% Requer \usepackage{xcolor} no preambulo para exibir os acrescimos em azul.
\section{Introduction}
\label{Introduction}

Os nematoides fitoparasitas constituem um dos principais fatores limitantes da produtividade agrícola, pois comprometem o desenvolvimento radicular e reduzem a capacidade das plantas de absorver água e nutrientes. Entre as espécies de maior importância para a cultura da soja destaca-se o nematoide de cisto, \textit{Heterodera glycines}, considerado um dos principais agentes de perdas produtivas dessa cultura \cite{bent2022}. A persistência de seus cistos no solo por períodos prolongados dificulta a erradicação do patógeno e reforça a necessidade de métodos capazes de identificar áreas infestadas e apoiar a definição de estratégias de manejo.

O diagnóstico convencional de \textit{H. glycines} baseia-se principalmente na coleta de amostras de solo e raízes, que são encaminhadas a laboratórios especializados para extração de cistos e ovos, contagem e identificação do nematoide. O processamento envolve etapas sucessivas de peneiramento, lavagem, separação e análise microscópica, cuja execução demanda tempo e profissionais capacitados \cite{arjoune2022,mahecha2022,pun2024}. Métodos moleculares e outras abordagens de diagnóstico também podem ser empregados, conforme a finalidade da análise e os recursos disponíveis \cite{shao2023}. \textcolor{blue}{A coleta, o acondicionamento e o transporte integram esse fluxo, mas seus custos e prazos dependem das condições locais e não foram medidos neste estudo. Assim, uma triagem portátil poderia futuramente apoiar a escolha de amostras para confirmação laboratorial; seu efeito sobre a logística ainda precisaria ser avaliado.}

Outro desafio está relacionado à distribuição espacial do nematoide no campo. Populações de \textit{H. glycines} podem apresentar padrões de agregação e diferenças de densidade dentro de um mesmo talhão \cite{avendano2003,poromarto2019}. Consequentemente, a representatividade do diagnóstico depende da quantidade e da distribuição dos pontos amostrados. \textcolor{blue}{Essa heterogeneidade motiva o estudo de ferramentas de triagem que possam, após validação em campo, auxiliar na seleção de pontos para análises nematológicas confirmatórias.}

Tecnologias baseadas em sensores têm sido investigadas como alternativas para o monitoramento das condições do solo e das plantas. Sensores de pH, condutividade elétrica, umidade e temperatura, assim como sistemas ópticos de sensoriamento remoto, permitem acompanhar variáveis relacionadas ao desenvolvimento das culturas \cite{santos2022}. Entretanto, essas variáveis apresentam especificidade limitada para a detecção de nematoides, uma vez que diferentes fatores bióticos e abióticos podem produzir respostas semelhantes. A análise de compostos orgânicos voláteis (VOCs) constitui uma abordagem complementar, pois alterações fisiológicas decorrentes da interação entre planta e patógeno podem modificar o perfil volátil do sistema biológico. Perfis de VOCs também respondem a estresses não relacionados ao patógeno investigado; por isso, uma diferença entre grupos experimentais não comprova, isoladamente, especificidade para \textit{H. glycines} \cite{gan2023}.

Alterações na emissão de VOCs de plantas de soja infestadas por \textit{H.~glycines} já foram investigadas por cromatografia gasosa acoplada à espectrometria de massas (GC--MS). \citet{constantino2021} observaram diferenças entre os perfis voláteis da parte aérea de plantas infestadas e não infestadas, com separação das condições experimentais a partir do quinto dia após a infestação, antes do aparecimento de diferenças significativas no porte das plantas. Esses resultados indicam que a interação entre a soja e o nematoide pode produzir alterações detectáveis na composição volátil. Contudo, o estudo concentrou-se nos compostos emitidos pela parte aérea, mantendo o solo e o sistema radicular isolados durante a coleta. Em outro contexto, \citet{velloso2021} mostraram que compostos liberados por diferentes plantas podem modificar a eclosão de ovos de \textit{H.~glycines}. Esse achado sustenta a relevância biológica dos voláteis na interação planta--nematoide, mas não demonstra um método de diagnóstico.

\textcolor{blue}{Permanece, portanto, uma lacuna experimental delimitada: avaliar se os gases coletados de vasos contendo solo e plantas de soja produzem sinais capazes de distinguir vasos alocados ao grupo inoculado de vasos controle em condições controladas. A hipótese deste trabalho é que a inoculação possa alterar, direta ou indiretamente, a composição da mistura gasosa amostrada do sistema solo--planta.} O padrão registrado pode refletir respostas fisiológicas da soja, mudanças na atividade microbiana do substrato, efeitos do material utilizado na inoculação ou alterações gerais do crescimento vegetal. A atividade microbiana do solo, por si, pode gerar sinais detectáveis por narizes eletrônicos \cite{decesare2011}. Assim, a discriminação das classes experimentais não permite identificar a origem química de cada sinal nem concluir que os sensores respondem diretamente ao nematoide.

Embora a GC--MS permita separar e identificar individualmente os compostos presentes em uma amostra, sua aplicação geralmente requer instrumentação especializada, preparação das amostras e operação por pessoal capacitado. O tamanho e as exigências operacionais dos equipamentos também limitam sua utilização em medições rápidas e diretamente em condições de campo \cite{banerji2024}. Em contraste, os sistemas de nariz eletrônico utilizam matrizes de sensores de gases com seletividade parcial, cujas respostas combinadas formam um padrão multivariado associado à composição volátil da amostra. Esse padrão, denominado assinatura olfativa, pode ser processado por técnicas de reconhecimento de padrões e aprendizado de máquina para diferenciar amostras sem a necessidade de identificar individualmente cada composto \cite{berna2010,li2024enose}. Embora não substituam a capacidade analítica da GC--MS, os narizes eletrônicos apresentam potencial para medições rápidas, portáteis e com menor complexidade operacional \cite{li2024enose}. Ao mesmo tempo, apresentam limitações de sensibilidade e especificidade, deriva de sinal e interferência das condições ambientais, razão pela qual devem ser avaliados como ferramentas complementares ao diagnóstico confirmatório \cite{cellini2017}.

\textcolor{blue}{Nesse contexto, três perguntas orientam a avaliação exploratória: (i) os sinais dos sensores de gases da matriz MQ diferenciam os grupos experimentais de vasos? (ii) a separação observada se mantém sob diferentes condições de umidade do solo e de aquisição? (iii) os sinais MQ acrescentam informação para a classificação além da leitura de umidade do solo, da temperatura e da pressão?} O teor de água do solo pode alterar a liberação e a retenção de voláteis, enquanto a umidade do gás e a temperatura podem interferir na resposta sensorial \cite{cellini2017}. \textcolor{blue}{A pressão registrada durante a aquisição descreve também o funcionamento da bomba e o transporte da amostra gasosa até a câmara sensorial. Por isso, umidade, temperatura e pressão podem contribuir para a previsão das classes sem representar, por si, um sinal específico do nematoide. A comparação entre modelos com sensores MQ, com variáveis ambientais e operacionais e com ambos os conjuntos permite examinar esse possível fator de confusão.}

\textcolor{blue}{Este estudo avaliou, em ambiente controlado, um nariz eletrônico aplicado aos gases coletados do sistema solo--soja para comparar vasos alocados ao grupo experimentalmente inoculado com \textit{H. glycines} e ao grupo controle. Investigou a contribuição dos sinais MQ em relação às leituras de umidade do solo, temperatura e pressão, considerando as condições de umidade e o protocolo de aquisição. Os rótulos descrevem a alocação experimental dos vasos; a classificação obtida não estima a densidade do nematoide nem demonstra a detecção de infestações naturais. A aplicação da abordagem à triagem em campo, assim como seus possíveis ganhos logísticos, depende de validação externa e de avaliação operacional posterior.}
